\documentclass[11pt,a4paper]{article}
\usepackage[T1]{fontenc}
\usepackage[utf8]{inputenc}
\usepackage{lmodern,amsmath,amssymb,amsthm,mathtools,booktabs}
\usepackage[margin=27mm]{geometry}
\usepackage[hidelinks]{hyperref}
\usepackage{microtype}
\newtheorem{theorem}{Theorem}[section]
\newtheorem{lemma}[theorem]{Lemma}
\newtheorem{proposition}[theorem]{Proposition}
\newtheorem{remark}[theorem]{Remark}
\newcommand{\R}{\mathbb R}
\newcommand{\Z}{\mathbb Z}
\newcommand{\im}{\operatorname{im}}
\newcommand{\rank}{\operatorname{rank}}
\newcommand{\norm}[1]{\lVert#1\rVert}
\newcommand{\ip}[2]{\langle#1,#2\rangle}
\DeclareMathOperator{\Rea}{Re}
\providecommand{\SourceCommit}{\errmessage{A frozen SourceCommit must be supplied}}
\pdfinfoomitdate=1
\pdftrailerid{}
\ifdefined\pdfsuppressptexinfo\pdfsuppressptexinfo=15\fi
\title{Canonical and unrestricted observations in a finite checkerboard gradient model}
\author{Lluis Eriksson\\\small THE ERIKSSON PROGRAMME}
\date{11 November 2026}
\begin{document}
\maketitle
\begin{abstract}
We compare two classes of linear observations for a finite, source-weighted,
two-channel checkerboard gradient form. For orthogonal alternating weak lines
in dimensions three and four, the minimum restoring unrestricted rank is
exactly $Dn$ on all fields and $(D-1)(n-2)$ on the actual coexact domain.
For any fixed positive angle, an explicit transverse Fourier cutoff gives
an unrestricted observation count of order $n$ and a quadratic penalty
independent of $n$, with constants depending on the angle and dimension.
A real coexact counterexample shows that the mean-only orthogonal construction
does not extend unchanged to small positive angles. In contrast, canonical
original-channel rectangle means have a necessary ambient count of order
$n^{D-2}$. This last argument does not apply to the coexact domain, because
rectangular averaging need not preserve it. A further explicit coexact witness
excludes partitions whose every cell in either channel leaves a coordinate
uncut. For a three-dimensional product-slice class, an explicit tagged-window
cutoff proves a necessary inequality involving corner matching over every
number of slices. Fixed or bounded-motion cuts require a logarithmic loss
in full-row count, and a rapid helical grid fails through a short resonance.
The general canonical coexact three-dimensional question, linear count
versus a possible logarithmic loss, remains open. These are finite analytical results;
we establish neither historical priority, formal verification nor a physical
Yang--Mills spectral gap.
\end{abstract}

\section{Problem and scope}
Spectral coarse spaces that retain modes responsible for slow convergence are
established constructions; see Spillane et al.~\cite{spillane2011,spillane2014}.
The calculations here concern a particular finite coefficient family and a
restoring quadratic inequality. They do not prove a preconditioned condition
number bound or a new general principle of spectral enrichment. Our purpose is
to keep the coefficients, observed rows and domain explicit enough to compare
their different rank and count requirements.

The manuscript consolidates the programme's analytical sources 251--254 and 256--259,
rather than presenting their corollaries as separate papers. A preceding
working paper treats broader canonical constructions~\cite{eriksson250}.
The present manuscript is a proposed substantial revision of the existing
checkerboard working paper~\cite{eriksson255}. Its new PDF is ungraded at
this source stage; earlier assessments apply only to their original bytes.

Fix an integer $D\geq3$, an even integer $n\geq2$, and $c>0$. Write
\[
 X=(\Z/n\Z)^D,\quad V=n^D,\quad
 a=\frac{c}{16D},\quad H=cn^2.
\]
We identify the original chart with $\{0,\ldots,n-1\}^D$ and retain all
positive periodic labels $\partial_i f(x)=f(x+e_i)-f(x)$, including at $n=2$.
The source parity is $p(x)=\sum_i x_i\pmod2$. For each parity $l=0,1$,
let $v_l,s_l\in\R^2$ be an orthonormal pair. The two weak lines
$\R v_0,\R v_1$ have principal angle $\theta\in(0,\pi/2]$. Put
\[
 B(x)=a v_{p(x)}v_{p(x)}^T+H s_{p(x)}s_{p(x)}^T.
\]
An original field is $z=(z_1,\ldots,z_D)$, with $z_r(x)\in\R^2$.
All inner products and norms use counting measure, summing sites, spatial
components and internal channels. Define
\begin{equation}\label{eq:form}
 Q(z)=\sum_{x,i,r}(\partial_i z_r(x))^T B(x)\partial_i z_r(x),
 \qquad E(z)=\sum_{x,i,r}|\partial_i z_r(x)|^2.
\end{equation}
The weights in \eqref{eq:form} belong to the source $x$.

For a real spatial one-form $u$, use the forward curl
$(d_1u)_{ij}=\partial_i u_j-\partial_j u_i$, with independent indices $i<j$.
Its adjoint is taken for counting measure on those independent components.
Set $U_0=\im(d_1^*)$ and let $U=U_0^2$ mean one copy for each original
internal channel. This is the actual coexact domain, not a substitute
divergence-free space or an abstract low-energy band.

For a domain $A$, a linear map $C:A\to F$ restores target $c$ if
\begin{equation}\label{eq:restore}
 V^{-1}Q(z)+K\norm{Cz}^2\geq cV^{-1}\norm{z}^2\quad(z\in A)
\end{equation}
for some finite $K\geq0$. The lower bounds below also apply to any finite
penalty $\Phi(Cz)$ with $\Phi(0)=0$. The rank on $U$ is always the actual
restricted rank $\rank(C|_U)$.

A canonical map uses a nonempty integer rectangle partition $\mathcal P_\alpha$
of the original chart for each original channel $\alpha=1,2$, shared by all
spatial components. Its rows and full count are
\begin{equation}\label{eq:canonical}
 (Cz)_{\alpha,r,S}=V^{-1/2}|S|^{-1/2}\sum_{x\in S}z_{r,\alpha}(x),
 \qquad N_{\rm out}=D(|\mathcal P_1|+|\mathcal P_2|).
\end{equation}
The rectangles are products of nonwrapping chart intervals. On the ambient
space these disjoint, nonzero rows are orthogonal and have rank $N_{\rm out}$.
On $U$, that equality is not presumed. Unrestricted maps may instead use
arbitrary real linear rows, including internal-channel mixtures and nonlocal
Fourier coefficients.

\section{Actual strong components and their spectrum}\label{sec:graph}
Let $S$ be the $2\times2$ matrix with rows $s_0^T,s_1^T$ and put
\begin{equation}\label{eq:delta}
 Y_{l,r}=s_l\cdot z_r,\qquad
 \delta=1-|s_0\cdot s_1|=1-\cos\theta>0.
\end{equation}
The row Gram matrix has eigenvalues $1\pm|s_0\cdot s_1|$. Thus $S$ is
invertible and $\norm{Sz}^2\geq\delta\norm{z}^2$. This is a proof coordinate;
we will specify its resulting observations as actual original-field rows.
Dropping only nonnegative weak terms gives
\begin{equation}\label{eq:strong}
 Q(z)\geq H E_{\rm str}(Y),\qquad
 E_{\rm str}(Y)=\sum_{l,r}\sum_{x:p(x)=l}\sum_i
 |Y_{l,r}(x+e_i)-Y_{l,r}(x)|^2.
\end{equation}

For each residue $s$ of parity $l$, define
\[
 G_{l,s}=\{x:\textstyle\sum_i x_i\in\{s,s+1\}\pmod n\}.
\]
For each $l$ these $n/2$ sets partition $X$. They include the wrap pair
$\{n-1,0\}$ where appropriate. Put $p=D-1$ and, for $y\in(\Z/n\Z)^p$, set
\[
 x_s(y)=(y_1,\ldots,y_p,s-\textstyle\sum_i y_i),\quad
 u(y)=Y_{l,r}(x_s(y)),\quad t(y)=Y_{l,r}(x_s(y)+e_D).
\]
These are bijections onto the actual source and destination layers. Since
$x_s(y)+e_i=x_s(y+e_i)+e_D$ for $i<D$, the component energy is exactly
\begin{equation}\label{eq:layers}
 E_s=\sum_y|u(y)-t(y)|^2+
       \sum_{y,i<D}|u(y)-t(y+e_i)|^2.
\end{equation}
Rungs and cross edges generate every source translation, and every target
has a rung. The sets $G_{l,s}$ are therefore connected actual strong
components. No labelled edge has been reassigned or deduplicated.

At a counting-orthonormal Fourier frequency $k\in(\Z/n\Z)^p$, put
$T_k=1+\sum_{i=1}^p e^{2\pi i k_i/n}$. The Hermitian symbol of
\eqref{eq:layers} has diagonal $D$ and off-diagonal entries $-T_k$ and its
conjugate. Its eigenvalues are $D-|T_k|$ and $D+|T_k|$. Expansion gives
\begin{equation}\label{eq:spectrum}
 D^2-|T_k|^2=\sum_{0\leq i<j\leq p}|e^{i\vartheta_i}-e^{i\vartheta_j}|^2,
 \quad \vartheta_0=0,\quad\vartheta_i=2\pi k_i/n.
\end{equation}
Consequently
\begin{equation}\label{eq:gap-symbol}
 D-|T_k|\geq\frac1{2D}\sum_{i=1}^p|1-e^{2\pi i k_i/n}|^2.
\end{equation}
At nonzero $k$, at least one squared term is at least
$4\sin^2(\pi/n)\geq16/n^2$. At zero frequency the eigenvalues are $0,2D$;
the zero vector branch is precisely the common constant on both layers.
The other branch is also at least $8/(Dn^2)$ for $D\geq3,n\geq2$.
Finite Parseval therefore proves
\begin{equation}\label{eq:mean-gap}
 E_s(f)\geq\frac8{Dn^2}\norm{f-\operatorname{mean}_{G_{l,s}}f}^2.
\end{equation}
Conjugate Fourier frequencies pair in the real space; real Nyquist
frequencies obey the same real matrix inequality. All formulas include $n=2$
with its original positive labels.

\section{Exact unrestricted ranks at the orthogonal angle}
Assume $D=3$ or $4$, and fix the original basis by
$v_0=e_1,v_1=e_2$, so channel 1 is strong on odd starts and channel 2 on even
starts. Let $R_{\rm str}$ average each original channel and spatial component
over its correct strong components. This is an ambient counting-orthogonal
projection. Define $C_{\rm str}$ by the normalized complete component means
as in \eqref{eq:canonical}, replacing rectangles by the corresponding
$G_{l,s}$. There are $Dn$ full rows and
\[
 \norm{C_{\rm str}z}^2=V^{-1}\norm{R_{\rm str}z}^2.
\]

\begin{theorem}\label{thm:exact}
For the orthogonal family in $D=3,4$, the minimum unrestricted restoring
ambient rank is exactly $Dn$. The minimum actual coexact restoring rank is
exactly $(D-1)(n-2)$, including zero at $n=2$. The same displayed map
$C_{\rm str}$ attains both ranks and restores target $c$ with $K=c$ on its
fixed normalized rows.
\end{theorem}
\begin{proof}
Sum \eqref{eq:mean-gap} over all original components. By \eqref{eq:strong},
\[
 Q(z)\geq\frac8D c\norm{(I-R_{\rm str})z}^2.
\]
Since $8/D\geq2>1$, orthogonality gives
$Q(z)+c\norm{R_{\rm str}z}^2\geq c\norm{z}^2$ on all original fields.
This proves the claimed normalized upper without any assertion about
preservation of $U$.

We next derive that preservation and the restricted rank for this particular
map. For one internal channel, write its sum-layer averages as
\[
 A_r(t)=n^{-(D-1)}\sum_{\sum x_i=t}z_r(x).
\]
A coexact field has zero componentwise mean and backward divergence zero,
by summing adjoint differences and using $d_1d_0=0$. Averaging the divergence
over layer $t$ shows that $F(t)=\sum_r A_r(t)$ satisfies $F(t-1)-F(t)=0$.
Its total sum is zero, so $F=0$. On a strong pair $\{s,s+1\}$ the projected
value is $[A_r(s)+A_r(s+1)]/2$. Hence $R_{\rm str}z$ depends only on
$t=\sum x_i$, has zero mean, has spatial sum zero pointwise, and is constant
on the appropriate strong pairs in each channel.

Conversely, any real mean-zero spatial vector function of $t$ whose spatial
sum is zero is actually coexact. At a nonzero diagonal frequency
$(j,\ldots,j)$ the forward multipliers are the common value
$\xi_i=e^{2\pi i j/n}-1$. For its coefficient $u$, $\sum_i u_i=0$, and hence
$\sum_i\overline{\xi_i}u_i=0$. The independent two-form coefficients
\begin{equation}\label{eq:preimage}
 \psi_{ik}=\frac{\xi_i u_k-\xi_k u_i}{\sum_i|\xi_i|^2},\quad i<k,
\end{equation}
have actual curl adjoint $u$, since expansion yields
\[
 (d_1^*d_1u)_r=(\textstyle\sum_i|\xi_i|^2)u_r
       -\xi_r\textstyle\sum_i\overline{\xi_i}u_i.
\]
There is no zero frequency; conjugation and the real Nyquist mode make the
inverse two-form real. Thus membership in $\im(d_1^*)$ has been proved.

Let $W_0$ be these functions, additionally constant on each channel's strong
pairs. With $m=n/2$, the scalar pair-constant mean-zero space has dimension
$m-1$. Choose $D-1$ such functions per channel and set the last to minus their
sum. It follows that
\[
 R_{\rm str}(U)=W_0,\quad \dim_\R W_0=2(D-1)(m-1)=(D-1)(n-2).
\]
Indeed the forward inclusion was derived above, $W_0\subset U$ by
\eqref{eq:preimage}, and $R_{\rm str}$ is the identity on $W_0$. The kernels
of $C_{\rm str}|_U$ and $R_{\rm str}|_U$ agree, so its actual rank is this
dimension. Its ambient rank is $Dn$, because all full rows are disjoint.

Finally the full strong kernel consists of constants on each connected
component and has dimension $Dn$. Every field $w$ in it has
$Q(w)=aE(w)\leq4Da\norm w^2=c\norm w^2/4$, by the isometry of forward
shifts. Its intersection with $U$ is exactly $W_0$, by the two preceding
directions. Any restoring linear map, with any finite penalty vanishing at
zero, must be injective on each corresponding low-energy subspace. A nonzero
kernel vector there contradicts the target. These give both sharp lower
ranks and complete the proof.
\end{proof}

For $n\geq4$ a set $G_{l,s}$ projects onto the full chart in every coordinate,
but has size $2n^{D-1}<n^D$. It cannot be a product rectangle: all its factors
would then be full. Thus the attaining rows are not canonical rectangle rows.
The $n=2$ component is the whole torus and is the stated harmless exception.
This special projection preserves $U$ by proof; no such property is implied
for other block-averaging maps.

\section{An explicit upper at any fixed positive angle}
The angle assumption in Theorem~\ref{thm:exact} can be removed for an
order-$n$ unrestricted upper by retaining transverse modes. We do not extend
its exact minima to arbitrary angles.

\begin{theorem}\label{thm:angle}
For all $D\geq3$ and $\theta\in(0,\pi/2]$, put
\[
 R=\left\lceil\sqrt{\frac{D}{4\delta}}\right\rceil,\qquad
 L=\min(n,2R+1),\qquad K=\frac{2c}{\delta}.
\]
There is an unrestricted real map with $2DnL^{D-1}$ displayed rows satisfying
\begin{equation}\label{eq:angleupper}
 V^{-1}Q(z)+K\norm{Cz}^2\geq2cV^{-1}\norm z^2
\end{equation}
on all ambient fields and hence on $U$. Both unrestricted minimum restoring
ranks have order $n$ as even $n\to\infty$ at fixed $D,\theta>0$.
\end{theorem}
\begin{proof}
On each actual two-layer normal-coordinate component from
Section~\ref{sec:graph}, let $P$ retain both layer coefficients at frequencies
\[
 \mathcal A_R=\{k:|k_i|_n\leq R\text{ for every }i<D\},
\]
where $|k_i|_n$ is cyclic distance to zero. There are exactly $L^{D-1}$ such
frequencies. Conjugate pairs produce sine and cosine real coordinates, and
each self-conjugate character is counted once. The real dimension is therefore
$2L^{D-1}$ per component. The direct sum over all $l,s,r$ is an orthogonal
projection in the $Y$ norm.

Outside the cube, some cyclic coordinate distance is at least $R+1$ and at
most $n/2$. Thus \eqref{eq:gap-symbol} and $\sin(\pi q/n)\geq2q/n$ give
the lower eigenvalue bound $\gamma=8(R+1)^2/(Dn^2)$. The upper branch is
at least $D\geq\gamma$, because a nonempty tail implies
$\gamma\leq2/D$. By Parseval,
\[
 E_{\rm str}(Y)\geq\gamma\norm{(I-P)Y}^2.
\]
If the cube is the whole frequency torus, $P=I$ and the zero-tail statement
is automatic; no nonexistent exterior mode is invoked. This includes $n=2$.

Choose a real orthonormal character basis on each retained layer and extend
by zero to the other components; denote the resulting vectors by $b_j$.
They are explicit counting-orthonormal vectors in $Y$ coordinates. Define
actual original-field rows by
\[
 (Cz)_j=V^{-1/2}\ip{b_j}{Y_{l,r}},\qquad Y_{l,r}=s_l\cdot z_r.
\]
There are $2L^{D-1}$ rows per component, $n/2$ components per parity,
two parities and $D$ spatial components. Consequently
$N_{\rm out}=2DnL^{D-1}$ and $\norm{Cz}^2=\norm{PSz}^2/V$ exactly.
Our cutoff yields $H\gamma\geq2c/\delta=K$. Therefore
\[
 Q(z)+VK\norm{Cz}^2\geq
 K\norm{(I-P)Sz}^2+K\norm{PSz}^2
 =K\norm{Sz}^2\geq2c\norm z^2.
\]
The same conclusion holds for a full cube using only the observed term.
Orthogonality is used in the $Y$ norm, followed by the proved singular-value
bound in \eqref{eq:delta}. We have not assumed that $P$ or $S^{-1}PS$
preserves $U$; the ambient inequality is restricted only after its proof.

For the lower ranks, let $\mathcal F_l$ consist of mean-zero functions on
$\Z_n$ constant on the strong pairs $\{l+2j,l+2j+1\}$.
Its dimension is $m-1$, $m=n/2$. Choose independently
$f_{l,r}\in\mathcal F_l$ for $r<D$, put
$f_{l,D}=-\sum_{r<D}f_{l,r}$, and define original fields
\[
 z_r(x)=S^{-1}\begin{pmatrix}f_{0,r}(\sum x_i)\\ f_{1,r}(\sum x_i)\end{pmatrix}.
\]
This is injective and has real dimension $(D-1)(n-2)$. Every original channel
has zero mean and pointwise spatial sum zero; the explicit real preimage
\eqref{eq:preimage} proves membership in $U$. All actual strong normal
differences vanish. Thus $Q=aE\leq c\norm z^2/4$ on this whole subspace.
Every restoring map must inject it. Dropping the mean and spatial-sum
conditions gives $Dn$ independent ambient parameters with the same energy
bound, and hence the ambient lower rank $Dn$.
Since $R,K$ are independent of $n$ at fixed $D,\theta$, these lower ranks and
the constructed upper count give order $n$. The coexact lower is zero at
$n=2$; no positive endpoint lower is asserted.
\end{proof}

The count constant and penalty diverge as $\theta\to0$. At $\theta=0$ the
matrix $S$ is singular, and the theorem does not follow by taking a limit.
Even at a fixed angle, $K$ is a coefficient on these fixed normalized rows;
it is not the condition number of the full operator. The observations are
nonlocal and can mix the original internal channels. No canonical upper has
been obtained from this construction.

\section{Why complete component means do not suffice at small angles}
\begin{proposition}\label{prop:smallangle}
Let $v_0=(1,0)$, $v_1=(\cos\theta,\sin\theta)$ and
$0<\theta\leq1/4$ radians. For every even $n\geq2$ and $D\geq3$, there is a
nonzero real $z\in U$ annihilated by all complete strong-component means,
in any constant original-channel or normal-coordinate direction, with
$Q(z)/(c\norm z^2)<17/24$. Any finite penalty vanishing at zero therefore
fails to restore target $c$ for the map consisting only of those means.
\end{proposition}
\begin{proof}
Put $w=(\cos(\theta/2),\sin(\theta/2))$ and let
$z_D(x)=w\cos(2\pi x_1/n)$, $z_r=0$ for $r\ne D$.
At its supported original frequencies $\pm e_1$, only the first forward
multiplier is nonzero and only the $D$th spatial field component is nonzero.
Thus $\sum_i\overline{\xi_i}u_i=0$, and formula~\eqref{eq:preimage}, now
with these actual multipliers, constructs its real coexact preimage. At
$n=2$ the coincident Nyquist mode and preimage are real.

On any sum layer $s$, parameterize $x=(y,s-\sum y)$. Its cosine sum is
\[
 n^{D-2}\sum_{y_1=0}^{n-1}\cos(2\pi y_1/n)=0.
\]
Hence every $G_{l,s}$
mean vanishes, in every constant internal direction and every spatial
component. Normalization and linear combinations of these means do not
change this conclusion.

The bisector satisfies
$|v_l\cdot w|^2=\cos^2(\theta/2)$ and
$|s_l\cdot w|^2=\sin^2(\theta/2)$ for both parities. Thus its source
quadratic coefficient is the same number
$\beta=a\cos^2(\theta/2)+cn^2\sin^2(\theta/2)$ at every site.
The periodic cosine is an exact positive-Laplacian eigenfunction, so
\[
 \frac{Q(z)}{\norm z^2}=4\sin^2(\pi/n)
 \bigl[a\cos^2(\theta/2)+cn^2\sin^2(\theta/2)\bigr].
\]
This retains the forward labels also at $n=2$. Using $\sin t\leq t$ and
$\sin^2(\pi/n)\leq1$ gives
\[
 \frac{Q(z)}{c\norm z^2}\leq\frac1{4D}+\pi^2\theta^2
 \leq\frac1{12}+\frac{\pi^2}{16}<\frac{17}{24}<1.
\]
The last strict bound uses $\pi^2<10$. Since every observation is zero,
the target fails for every admitted penalty. The strong energy is positive
at positive angle; it was not incorrectly set to zero.
\end{proof}

This proposition refutes the unchanged mean-only extension for a specific
angle interval. It neither locates a sharp critical angle nor excludes the
enriched observations in Theorem~\ref{thm:angle}, which retain transverse
modes missed by these means.

\section{A canonical ambient count lower bound}\label{sec:canonical}
Return to the orthogonal original basis $v_0=e_1,v_1=e_2$.
Let $R$ be the canonical counting-orthogonal block projection, so
$\ker C=\ker R$ and $\norm{Cz}^2=\norm{Rz}^2/V$.

\begin{theorem}\label{thm:ambientcanonical}
If a canonical map restores target $c$ on all ambient fields, with any finite
penalty $\Phi(Cz)$ satisfying $\Phi(0)=0$, then
\begin{equation}\label{eq:countlower}
 N_{\rm out}\geq
 \frac{3DV}{2^{2D+4}(1+8Dn^2)}
 \geq\frac{3D}{2^{2D+4}(1+8D)}n^{D-2}.
\end{equation}
The statement concerns the ambient count and rank. It does not assert this
lower bound on the actual coexact domain alone.
\end{theorem}
\begin{proof}
First suppose $n\geq4$. Set $m=n/2$, $k=\lfloor m/2\rfloor$ and
$\varphi=2\pi k/n\in[\pi/3,\pi/2]$. On sum residues put
$h(2j)=h(2j+1)=\cos(2\varphi j)$. Choose a unit spatial vector $b$ and
set $z_{r,2}(x)=b_r h(\sum_i x_i)$, $z_{r,1}=0$. Every even-source strong
difference of channel 2 vanishes. Finite character sums give
\[
 \norm z^2=2n^{D-1}\sum_j\cos^2(2\varphi j)\geq V/2.
\]
The cyclic labelled difference identity on the $m$ pair values yields
\[
 Q(z)=2aD\sin^2\varphi\norm z^2\leq(c/8)\norm z^2.
\]
The period-two pair cycle retains both forward labels, so this includes $n=4$.

Write $A=(1+e^{-i\varphi})/2$, $B=(1-e^{-i\varphi})/2$.
The paired complex carrier is
$f(t)=Ae^{i\varphi t}+Be^{i(\varphi+\pi)t}$ and $h=\Rea f$.
Its two frequencies are well-defined on $\Z_n$. Each interval character sum
is at most $2/|1-e^{i\omega}|\leq2$ for these frequencies, and
$|A|+|B|\leq\sqrt2$. Factoring over every original rectangle $S$ gives
\[
 \left|\sum_{x\in S}h(\sum_i x_i)\right|^2\leq2^{2D+1}.
\]
For $w=z/\norm z$, put $t=\norm{Rw}^2$. With the shared spatial-component
partition and nonempty cells,
\[
 t\leq\frac{2^{2D+2}}V\sum_{S\in\mathcal P_2}\frac1{|S|}
 \leq\frac{2^{2D+2}N_{\rm out}}{DV}.
\]
The actual field $u=(I-R)w$ has $Cu=0$ and $\norm u^2=1-t$.
Ambient restoration requires $Q(u)\geq c(1-t)$.
On all ambient fields, $B\leq HI$ and forward-shift isometry imply
$Q(y)\leq4DH\norm y^2$. Positivity of the form then gives
\[
 Q(w-Rw)\leq2Q(w)+2Q(Rw)\leq c/4+8Dcn^2t.
\]
Consequently $t\geq3/[4(1+8Dn^2)]$. Combining these bounds proves the
first inequality in \eqref{eq:countlower}; the second follows from
$1+8Dn^2\leq(1+8D)n^2$.

At $n=2$, use no nonexistent oscillatory pair mode. Each channel partition
has at least one cell, so $N_{\rm out}\geq2D$. The first right-hand side of
\eqref{eq:countlower}, divided by $2D$, equals
$3/[2^{D+5}(1+32D)]<1$. This proves the endpoint directly.
\end{proof}

In $D=4$ the second constant is $1/11264$. The earlier ambient canonical
construction in programme source 247, included in the public source package
of~\cite{eriksson250}, has count at most $2^{25}n^2$. Together these give
order $n^2$ for the ambient canonical class. Theorem~\ref{thm:exact} gives
the same-profile ambient unrestricted minimum $4n$. These are comparisons
of row classes on the same domain, not comparisons of unlike domains.
The earlier upper is a cited proved input, not a result rederived here.

The indispensable domain boundary in the lower proof is $u=(I-R)w$.
Even if $b$ is chosen with spatial sum zero so that $w$ is coexact, a
rectangular block projection need not preserve $U$. Thus $u$ has not been
shown coexact, and applying a coexact target inequality to it would be
invalid. The argument supplies no canonical coexact $n^2$ lower in $D=4$.
In $D=3$ the canonical coexact question of $O(n)$ versus a logarithmic loss
also remains open; Theorem~\ref{thm:angle} does not settle it.

\section{A coexact falsifier for uncut rectangular cells}
\begin{proposition}\label{prop:uncut}
In the orthogonal family, let $n\geq4$ be even. If every cell of the
canonical partition of either original internal channel has at least one
full coordinate interval, that map cannot restore target $c$ on $U$ with
any finite penalty vanishing at zero. The full coordinate direction may
differ from cell to cell; the other channel may even observe all sites.
\end{proposition}
\begin{proof}
Choose the indicated channel $\alpha$ and its strong-source parity
$p_1=1,p_2=0$. Set $m=n/2\geq2$ and
\[
 F_\alpha(t)=\cos\left(\frac{2\pi}{m}
 \left\lfloor\frac{(t-p_\alpha)\bmod n}{2}\right\rfloor\right).
\]
It is constant on each correct strong pair, including the wrap pair, has
zero sum and is nonzero. Put
$z_{r,\alpha}(x)=b_rF_\alpha(\sum_i x_i)$ for
$b=(1,-1,0,\ldots,0)$, and set the other channel to zero.
The mean is zero and the spatial sum is zero. Formula~\eqref{eq:preimage}
constructs its actual real coexact preimage. Its strong differences vanish,
so $Q(z)=aE(z)\leq c\norm z^2/4$.

For any cell in that channel, fix all coordinates except a complete one.
As the complete coordinate varies, the sum residue runs through every
element of $\Z_n$ once. Its $F_\alpha$ sum is zero; summing the remaining
cell coordinates preserves zero. Thus every spatial mean in that cell
vanishes. The same field works simultaneously for all cells, irrespective
of which coordinate is full. The other channel is identically zero.
Therefore $Cz=0$ and the nonzero coexact field contradicts the target.
\end{proof}

A necessary condition is consequently that each original channel have at
least one cell proper in every coordinate. This is not sufficient, gives
no asymptotic logarithmic count lower, and does not exclude arbitrary
$O(n)$ partitions with all coordinates cut. At $n=2$ the scalar zero-mean
one-pair space vanishes; the proposition deliberately excludes that case.

\section{Temporal matching of product slice corners}\label{sec:temporal}
We now restrict to $D=3$, $4\mid n$, $n\geq32$, and the orthogonal
internal basis. Thus $a=c/48$ and $H=cn^2$. This section concerns the actual
coexact domain $U=U_0^2$ and the original normalized means. It supplies a
count obstruction for a specified product-slice class, rather than an
unrestricted lower bound for every canonical partition.

At each longitudinal slice $x_1=i$, let $\{J_{i,j}\}_j$ and
$\{K_{i,k}\}_k$ be nonempty cyclic interval partitions of $\Z_n$ whose
product of factor counts is the same integer $q$ for every $i$. The cells
of channel1 are $\{i\}\times J_{i,j}\times K_{i,k}$. Their full spatial-row
count is $3nq$. The second channel can use any permitted partition. Write
$E_i$ for the product of the terminal endpoints of the two interval
partitions, so that $E_i$ has exactly $q$ distinct transverse points.
For $K\geq1$, let $S_K$ be the maximum over cyclic $i$ of the minimum
bottleneck matching distance between $E_i$ and $E_{i+K}$, using the
periodic transverse $\ell^1$ metric. These extrema exist over finite sets.
We impose no common labelling around the cycle.

\begin{theorem}\label{thm:temporal}
Put $L=\log((n/4+1)/3)$. If for some integer $K\geq1$,
\begin{equation}\label{eq:windowthreshold}
 L\geq 2^{18}q A_K,\qquad
 A_K=\max\{16K,9K+6S_K+S_K^2\},
\end{equation}
there exists a real nonzero $z\in U$ with every original mean zero and
$Q(z)<7c\norm z^2/16$. In particular, no finite quadratic penalty on those
means can restore target $c$. Consequently recovery requires
\begin{equation}\label{eq:windownecessary}
 L<2^{18}q\max\{16K,9K+6S_K+S_K^2\}
 \quad\hbox{for every integer }K\geq1.
\end{equation}
The assertion permits $K\geq n$ and includes the longitudinal seam.
\end{theorem}

We give the constructive proof, including the norm and original-source
energy bounds. No assumed capacity estimate, projection invariance or
coexact gap is used as an input.

\subsection{A finite transverse cutoff}
Let $Y=\Z_n^2$, $R=n/4$, and let $\rho_p$ be periodic $\ell^\infty$
distance from $p\in Y$. Set
\[
 f_p(y)=\min\{1,\max\{0,L^{-1}\log((\rho_p(y)+1)/3)\}\}.
\]
It is zero for $\rho_p\leq2$ and one for $\rho_p\geq R$. We claim
\begin{equation}\label{eq:radial}
 \sum_{y,h=2,3}|\partial_h f_p(y)|^2\leq\frac{64}{L},
 \qquad \norm{1-f_p}_Y^2\leq\frac{4n^2}{L^2}.
\end{equation}
For a nonconstant edge, the smaller radius is $2\leq r<R$. There are
$8r+4\leq32r$ forward-labelled edges crossing that square boundary,
disregarding direction. Since $r+1\leq R\leq n/4$, no antipodal ambiguity
enters this count. Each scalar increment is at most $1/((r+1)L)$.
Thus the energy is at most
\[
 \frac{32}{L^2}\sum_{r=2}^{R-1}\frac{r}{(r+1)^2}
 \leq\frac{32\log(R+1)}{L^2}\leq\frac{64}{L},
\]
where $R\geq8$ gives $\log(R+1)\leq2L$.
For the norm, group sites by
$(R+1)/2^{j+1}<\rho_p+1\leq(R+1)/2^j$, $j\geq0$.
Each group has at most $4(R+1)^2/4^j$ sites and deficit at most
$(j+1)\log2/L$. The latter estimate holds also where clipping makes
$f_p=0$. Since $\sum_{j\geq0}(j+1)^2/4^j=80/27<4$ and $R+1\leq n/2$,
the deficit norm is at most $16(R+1)^2/L^2\leq4n^2/L^2$.
This proves~\eqref{eq:radial}. Telescoping a shortest lattice path also gives
\begin{equation}\label{eq:singletranslation}
 \norm{f_p(\cdot+u)-f_{p'}}_Y
 \leq (|u|_1+d(p,p'))\sqrt{64/L}.
\end{equation}
Here $|u|_1$ can be any chosen integer representative length, and all seams
are periodic.

\subsection{Tagged windows and their translation cost}
Fix $K$ and write $S=S_K$. At slice $i$ take the tagged list
$F_i=\{(j,p):0\leq j<K,\ p\in E_{i+j}\}$. Repeated spatial points keep
their tags, so its size is exactly $Kq$ even if $K\geq n$.
Match $F_i$ to $F_{i+1}$ by keeping the designated $(K-1)q$ shared entries
at the same spatial position and matching the remaining $q$ outgoing
entries to $E_{i+K}$ with distance at most $S$. Some replaced entries can
also be stationary; no claim about the total number of stationary matches
is needed. Define
\[
 \eta(i,y)=\min_{(j,p)\in F_i} f_p(y).
\]
It lies in $[0,1]$ and vanishes within radius2 of $E_i$. Squared maximum
is bounded by the sum of squares, so~\eqref{eq:radial} gives
\begin{equation}\label{eq:etadeficit}
 \norm{1-\eta}_X^2\leq 4KqV/L^2.
\end{equation}
For two equally indexed real lists, $|\min a-\min b|\leq\max|a-b|$.
Using the designated matching and~\eqref{eq:singletranslation},
\begin{align}\label{eq:movingtranslation}
 \norm{\eta(i+1,\cdot+u)-\eta(i,\cdot)}_Y^2
 &\leq\frac{64q}{L}\{(K-1)|u|_1^2+(|u|_1+S)^2\}\notag\\
 &=\frac{64q}{L}(K|u|_1^2+2|u|_1S+S^2).
\end{align}
With no longitudinal shift, identity matching of all $Kq$ entries instead
gives $64Kq|u|_1^2/L$. Summing either bound over $i$ adds a factor $n$.
The window is only a cutoff recipe; the observed cells remain the original
$q$ cells per slice.

\subsection{Actual coexact membership and observation zeros}
Put $\delta_h f(x)=f(x)-f(x-e_h)$. For a real periodic $\Psi$, take
$\phi=\delta_2\delta_3\Psi$ and the actual two-form potential
$\psi_{23}=\phi$, $\psi_{12}=\psi_{13}=0$. Since
$\partial_h^*=-\delta_h$, direct adjunction yields
\[
 d_1^*\psi=(0,\delta_2\delta_3^2\Psi,-\delta_2^2\delta_3\Psi)=P\Psi.
\]
Every component has zero full sum by telescoping. Summing the second
component on a cell telescopes to four corner values of $\delta_3\Psi$;
the third gives the negative four corner values of $\delta_2\Psi$.
Predecessors of initial endpoints are terminal endpoints of the preceding
interval. Hence all relevant corners are in $E_i$, including wrap cells.
A whole-cycle factor cancels identically.

Use the period-four wave
$\Psi_0(x)=\sin(\pi(x_1+x_2+x_3)/2)$, and set
$\Psi=\eta\Psi_0$. At every corner, $\eta$ and its backward transverse
neighbors vanish. Thus both required differences of $\Psi$ vanish there.
Every channel1 mean of $P\Psi$ is exactly zero, with its original positive
normalization. Take channel2 identically zero. This constructs an actual
real $z=(P\Psi,0)\in U$ annihilating all observed rows.

\subsection{Nonzero norm and original energy}
Direct period-four differences give $z_0=P\Psi_0=(0,w,-w)$ with successive
$w$ values $(2,-2,-2,2)$. Thus $|z_0(x)|^2=8$,
$\norm{z_0}^2=8V$, and at every original odd source all $\partial_hz_0$
vanish, including across seams because $4\mid n$.
For a finite convolution $T=\sum_r a_r T^r$, subtraction gives
\[
 T(\eta\Psi_0)-\eta T\Psi_0
 =\sum_r a_r(T^r\eta-\eta)T^r\Psi_0.
\]
Each nonzero component of $P$ has coefficient absolute sum at most8,
longitudinal shift0 and transverse radius at most3. Since $|\Psi_0|\leq1$,
the no-longitudinal translation bound yields
\begin{equation}\label{eq:pnormerror}
 \norm{P(\eta\Psi_0)-\eta z_0}^2\leq73728nKq/L.
\end{equation}
Each component of $\partial_2P$ or $\partial_3P$ has coefficient absolute
sum at most16 and transverse radius at most4. Its squared translation
factor is at most $16K$. Each component of $\partial_1P$ has coefficient
absolute sum at most16, longitudinal shifts0 or1 and transverse radius
at most3. Its factors are bounded by $9K$ and $9K+6S+S^2$ respectively.
Consequently each derivative commutator has norm at most
$16\sqrt{64nqA_K/L}$. Restricting to original odd sources where
$\eta\partial_hz_0=0$, and summing the six squared component bounds,
\begin{equation}\label{eq:strongerror}
 \sum_{x\text{ odd},h}|\partial_h(P\Psi)(x)|^2
 \leq98304nqA_K/L.
\end{equation}
These are the original strong sources; no weight has been shifted.

By~\eqref{eq:etadeficit}, \eqref{eq:pnormerror} and $n\leq V$,
\[
 \norm{P\Psi-z_0}^2\leq
 \left(64Kq/L^2+147456Kq/L\right)V.
\]
The threshold and $A_K\geq16K$ imply $L\geq2^{22}Kq$. The coefficient
above is at most $2^{-38}+9/256<2$. Reverse triangle therefore gives
$\norm{P\Psi}>\sqrt{8V}-\sqrt{2V}=\sqrt{2V}$.
In particular the constructed field, not just its potential, is nonzero.
Multiplying~\eqref{eq:strongerror} by $H=cn^2$ gives
\[
 Q_{\rm strong}(z)\leq98304cVqA_K/L
 \leq3cV/8<3c\norm z^2/16.
\]
Every directional forward difference has norm at most $2\norm z$, so
$Q_{\rm weak}(z)\leq12a\norm z^2=c\norm z^2/4$.
Adding proves Theorem~\ref{thm:temporal}. Any quadratic penalty vanishes
on the identically zero observation vector.

\subsection{Count consequences and a rapidly moving grid}
For fixed cuts, $S_1=0$ and $A_1=16$. With bounded one-step corner
motion, $A_1$ stays bounded as well. In either class,
\eqref{eq:windownecessary} forces $q$ to grow at least proportionally to
$L$, and hence requires order $n\log n$ full channel1 rows. This conclusion
is restricted to these classes; arbitrary original-channel rectangle
partitions and rapidly moving product partitions remain outside it.

Large adjacent motion does not itself guarantee recovery. For
$n=2^m$, $m\geq5$, let $\kappa=n/4+2$ and $t_i=\kappa i\bmod n$.
Use the singleton $\{t_i\}$ and its cyclic complement for the first
transverse factor, and $\{0\}$ and its complement for the second.
Use these partitions in both original channels. There are $24n$ full rows
and $E_i=\{t_i,t_i-1\}\times\{0,n-1\}$, so $q=4$.
The identity $4\kappa=n+8$ gives $S_4\leq8$ and $A_4\leq148$.
For $L\geq2^{20}\cdot148$, the theorem gives an invisible real nonzero
coexact field below target. Such dyadic sizes exist, but the threshold is
enormous and is not a practical crossover estimate.

The adjacent bottleneck is exactly $\kappa$: the next corner with first
transverse coordinate $t_i+\kappa$ is at least that far from each old
corner, since $\kappa<n/2-1$; translation realizes that maximum distance.
Also $\gcd(\kappa,n)=2$, so the orbit visits $n/2$ positions.
This is a genuinely rapidly moving family with a short resonance. The
theorem does not assert that every rapidly moving $O(n)$ family fails.


\section{Interpretation and remaining question}
The sharp orthogonal ranks count the actual strong kernel and its explicitly
derived coexact image. Positive-angle enrichment pays a deteriorating
singular-value comparison by retaining finitely many transverse modes per
component. The temporal matching theorem supplies an actual invisible coexact field
for a specified product-slice class, including a rapidly moving resonant
example. Its restricted quantifiers do not settle the general canonical
count optimum. The small-angle counterexample explains why that enrichment cannot
simply be replaced by an assertion that the old means still suffice.
The canonical ambient lower instead uses a real oscillatory carrier and an
actual ambient kernel projection. Its failure to transfer to $U$ is an
explicit domain obstruction, not a technical detail that can be omitted.

The remaining question is whether, for the orthogonal $D=3$ family and actual
$U$, canonical original-channel integer rectangle partitions admit $O(n)$ full
count with some finite size-dependent quadratic penalty, or whether the
known logarithmic upper loss is unavoidable. The mean supports and Fourier
rows used in the unrestricted theorems are not permitted substitutes.
No theorem here establishes a size-uniform canonical penalty, efficient
implementation or a full-operator condition number. All results concern a
finite toy gradient family. The continuum and physical Yang--Mills mass gap
remain open; the programme's PR29 accounting remains $20/41$, with
\texttt{TermSource=0}.

\section{Model assessments and review context}
The analytical source texts251--254 and256--259 received complete involved
AI source readings. Their exact source bytes, full reports, objections and
context must remain in the reproducibility package. The temporal proof
uses the designated shared-slot wording specified in the retained259
clarification. These readings are not human peer reviews, blind readings,
independent provider samples or Lean verification. The exposed source
reviewer identity was Codex, based on the OpenAI GPT-6 family; exact served
backend, version and reasoning effort were not exposed. Reviewers shared
account, infrastructure and programme context and had prior involvement.

At this source stage the proposed revision has no exact-PDF numerical
assessment, successful production claim or visual-review claim. No grade is
assigned or inherited. The existing255B version1 PDF has SHA256
\begin{center}
\small\texttt{d45a17a4142204628189441b9133aa1027b0d7577a56dc5de6c576d2af6102e7}
\end{center}
Its first exact-PDF assessment recorded importance $4.00/10$, overall
$4/10$ stars, recommendation \texttt{major\_revision} and three material
objections open. Those are historical facts about that version, not
predetermined grades or objection states for a successor. Its full original
report must be retained unchanged. A new version requires its genuine first
assessment of the exact new PDF; repeated samples to improve a grade are
not evidence of scientific progress.

Analytical correctness is separate from originality and relevance. Spectral
enrichment is established literature. The finite radial cutoff, minimum
inequality and window-cover mechanisms are elementary; historical priority
of this particular application remains unestablished. The huge sufficient
threshold is not a demonstrated practical regime. Neither a deposit nor
Working-paper status constitutes editorial acceptance or human peer review.

\section*{Reproducibility and provenance}
Scientific inputs are frozen programme sources251--254 and256--259.
The source archive must retain their raw bytes, complete analytical reports
and context, this manuscript's mapping, the frozen Git commit and all actual
production logs. In particular the temporal theorem is derived in source259
at commit\texttt{2f370a7da850b8066a9afe95c523e49bd1ef725e}; its full audit
and clarification were integrated unchanged through merge
\texttt{8485cfb2b8b11bb5acd240d8a0b133457bbfe2a1}.
The separate golden-orbit working derivation260 is excluded from this
manuscript: no unreviewed frontier result is promoted here.

Linux compilation, extraction and page rendering are production checks;
they do not certify mathematics or transfer a model score. No such build
or exact-PDF review is claimed at this draft-source stage. Subsequent
production and review facts belong to dated sidecars linked to the exact
PDF hash. The original published version remains immutable. The frozen
production source is \texttt{\SourceCommit}.

\begin{thebibliography}{9}
\bibitem{spillane2011}
N. Spillane, V. Dolean, P. Hauret, F. Nataf, C. Pechstein and R. Scheichl,
A robust two-level domain decomposition preconditioner for systems of PDEs,
\emph{C. R. Acad. Sci. Paris, Ser. I} \textbf{349} (2011), 1255--1259.
\href{https://doi.org/10.1016/j.crma.2011.10.021}{doi:10.1016/j.crma.2011.10.021}.
\bibitem{spillane2014}
N. Spillane, V. Dolean, P. Hauret, F. Nataf, C. Pechstein and R. Scheichl,
Abstract robust coarse spaces for systems of PDEs via generalized eigenproblems
in the overlaps, \emph{Numer. Math.} \textbf{126} (2014), 741--770.
\href{https://doi.org/10.1007/s00211-013-0576-y}{doi:10.1007/s00211-013-0576-y}.
\bibitem{eriksson250}
L. Eriksson, Canonical rectangle observations for matrix gradients: fixed weak
spaces and checkerboard counterexamples, AIRR Working paper,
ARR-2026-5XANPRA7E891Q9FZ, version 1 (2026).
\url{https://airr.science/papers/ARR-2026-5XANPRA7E891Q9FZ/versions/v1/}.
\bibitem{eriksson255}
L. Eriksson, Canonical and unrestricted observations in a finite checkerboard
gradient model, AIRR Working paper, ARR-2026-597T88D7WZ98HR0Z,
version1 (2026).
\url{https://airr.science/papers/ARR-2026-597T88D7WZ98HR0Z/versions/v1/}.
\end{thebibliography}
\end{document}
